\section*{Chapter \ref{ch:rs:price-and-volume-features} --- Price and Volume Features}

\begin{solution}{exo:rs:price-and-volume-features:1}
$\ln(51.20/49.80) = 0.0277$; $(0.0277)^2/(4\ln2) = 0.000277$, a daily volatility of 1.67\%.
\end{solution}

\begin{solution}{exo:rs:price-and-volume-features:2}
3-1 momentum skips the last month: $1.02 \times 0.99 \times 1.04 - 1 = 5.02\%$. The four-month return is $1.050 \times 1.03 - 1 = 8.17\%$.
\end{solution}

\begin{solution}{exo:rs:price-and-volume-features:3}
$1.2\%/40 = 0.03\%$ per USD~million; $2.5\%/2 = 1.25\%$ per USD~million, about 40 times less liquid by this measure.
\end{solution}

\begin{solution}{exo:rs:price-and-volume-features:4}
$w_j = (5 - j)/6 - (2 - j)^+/3$: $w_0 = 1/6$, $w_1 = 1/3$, $w_2 = 1/2$, $w_3 = 1/3$, $w_4 = 1/6$; the sum is $1.5 = (6 - 3)/2$.
\end{solution}

\begin{solution}{exo:rs:price-and-volume-features:5}
The variance of an average falls with the number of days, so $60/5.2 = 11.5$ days of Parkinson estimates match sixty squared
returns.
\end{solution}

\begin{solution}{exo:rs:price-and-volume-features:6}
$(\sqrt{11} \times 0.03 + 1 \times (-0.03))/\sqrt{12} = (0.0995 - 0.03)/3.464 = 0.020$, a third less than the 0.03 of the 12-1 part
alone: with a monthly reversal, the skip is worth a third of the IC.
\end{solution}

\begin{solution}{exo:rs:price-and-volume-features:7}
With 30 steps the range estimators read 22--30\% low (Parkinson 0.78, Garman--Klass and Rogers--Satchell 0.70) against 8--10\% with 390:
the true high and low fall between the sampled points more often and by more when the points are sparser. \omterm{def:rs:market-data-for-research:bar}{Bars} built from few trades
(illiquid stocks) have the same problem.
\end{solution}

\begin{solution}{exo:rs:price-and-volume-features:8}
A window centred on today includes the next ten days' volume, which is not known at the close: the feature looks ahead. The leakage
test would catch it. Use the preceding 21 days, excluding today.
\end{solution}

\begin{solution}{pb:rs:price-and-volume-features:1}
\textbf{1.} For a Brownian motion observed continuously, the expected squared range of its logarithm over a day is $4\ln2\,\sigma^2$.
\textbf{2.} Parkinson. \textbf{3.} Rogers--Satchell: each of its two products pairs a high (or low) measured from the close and from
the open, so a drift enters with opposite signs and cancels. \textbf{4.} The overnight variance; $k = 0.34/(1.34 + 21/19) = 0.139$.
\textbf{5.} 5.2, 7.6 and 5.7. \textbf{6.} 7.5. \textbf{7.} About 7.6. \textbf{8.} 0.92, 0.90, 0.90 and 0.91: the range seen at 390
points is narrower than the continuous one. \textbf{9.} Close-to-close rises to 1.25 (it measures $\sigma^2 + \mu^2$), Parkinson to
1.02, Rogers--Satchell stays at 0.90. \textbf{10.} The single-day range estimators fall to 0.63--0.64; Yang--Zhang stays at 0.93.
\textbf{11.} The close-to-close return spans the night, so it contains the gap, and the drift adds $\mu^2$ to its second moment.
\textbf{12.} Highs at the ask and lows at the bid widen the range and bias every range estimator upwards, more for liquid-looking but
wide-spread stocks. \textbf{13.} Yang--Zhang. \textbf{14.} Rogers--Satchell (or Yang--Zhang with its overnight part near zero).
\textbf{15.} Close-to-close (an EWMA of squared returns), with five to eight times more days for the same precision. \textbf{16.} That
highs and lows exclude erroneous prints and auction prints, that the \omterm{def:rs:market-data-for-research:bar}{bars} cover the whole session, and how wide the spread is.
\textbf{17.} A less noisy volatility makes a better volatility \omterm{def:rs:anatomy-of-a-predictor:predictor}{predictor} and better risk scaling of other \omterm{def:rs:anatomy-of-a-predictor:predictor}{predictors}. \textbf{18.}
\emph{Named result:} relative to squared close-to-close returns, Parkinson is 5.2 times as efficient, Garman--Klass 7.6 and
Rogers--Satchell 5.7, and twenty-day Yang--Zhang 7.5 times the twenty-day close-to-close variance; with an overnight gap carrying 30\%
of the variance the single-day range estimators see 63--64\% of it. \textbf{19.} Low volatility and \omterm{def:rs:price-and-volume-features:volume}{Amihud illiquidity} (through absolute
returns), and any feature scaled by volatility. \textbf{20.} The path: how far the price travelled during the day, not only where it
ended.
\end{solution}

\begin{solution}{iq:rs:price-and-volume-features:1}
Because the most recent month's return tends to reverse (a negative first-order autocorrelation of monthly returns), so it works against
the continuation that momentum bets on; skipping it removes a component with the wrong sign. In a market without that reversal the skip
would cost information.
\iqlookfor{the reversal as the reason, not ``convention''.}
\end{solution}

\begin{solution}{iq:rs:price-and-volume-features:2}
Use the range: Garman--Klass or, with overnight gaps, Yang--Zhang, over the month, after cleaning highs and lows of bad prints; each is five
to eight times more efficient than squared close-to-close returns, and Yang--Zhang handles drift and gaps. Correct for the downward bias
of a discretely observed range if the \omterm{def:rs:market-data-for-research:bar}{bars} are thin.
\iqlookfor{range estimators, their efficiency and their failure cases.}
\end{solution}

\begin{solution}{iq:rs:price-and-volume-features:3}
The log of an average of prices is, to first order, an average of log prices, and each log price is today's minus a sum of recent returns.
The difference of the two averages gives $w_j = (m - 1 - j)/m - (n - 1 - j)^+/n$: triangular weights rising over the short window and
falling to zero at the long one.
\iqlookfor{the first-order derivation and the shape of the weights.}
\end{solution}

\begin{solution}{iq:rs:price-and-volume-features:4}
A perturbation test: recompute every feature after multiplying all inputs after a cut-off by random factors, and require identical outputs
up to the cut-off, for many cut-offs; plus a feature that deliberately peeks, to prove the test catches it. Also a data-access layer that
refuses data after the decision time.
\iqlookfor{an automated, falsifiable test.}
\end{solution}

\begin{solution}{iq:rs:price-and-volume-features:5}
The price move per unit of money traded, a proxy for price impact. Weaknesses: it is noisy day by day (a single large move dominates), mixes
news-driven moves with impact, depends on price level through the traded value, and is zero on days without trading.
\iqlookfor{what it proxies and why it is noisy.}
\end{solution}

\begin{solution}{iq:rs:price-and-volume-features:6}
The feature may carry real information that the simulator does not contain: a simulator only reproduces the effects planted in it. Nothing
is learned about the feature from the simulator; the real-data result must be tested like any other (search, overlap, costs), and the
simulator's scope is clarified: it cannot validate that feature's effect.
\iqlookfor{a simulator's scope, and not treating it as evidence against a real effect.}
\end{solution}
